\section{问题重述}

\subsection{问题背景}

水泥工业是国民经济的基础产业，同时也是高能耗、高排放的重点管控行业。在新型干法水泥生产流程中，窑尾熟料烧成环节是粉尘污染物的主要排放源。电除尘器（ESP）因其处理风量大、除尘效率高、运行阻力小等显著优势，已成为水泥生产线窑尾烟气净化的核心设备，其运行效果直接关系到企业的环保达标与能耗控制水平。

电除尘器利用高压电场使含尘气体中的粉尘荷电，并在电场力作用下将其从气流中分离捕集。在实际生产中，其除尘效率与运行能耗受到烟气温度、流量、入口浓度等工况参数，以及二次电压、振打周期等操作参数的多重耦合影响。当前，多数企业仍依赖人工经验设定参数，难以应对工况的剧烈波动，极易陷入“过度除尘、电耗居高不下”或“参数偏低、粉尘排放超标”的两难困境。面对动态变化的生产环境，传统手动调节方式响应滞后、精度不足，无法实现环保达标、能耗最低与稳定运行的协同目标。

因此，亟需构建一种智能化、精细化的协同优化控制模型。本问题基于某$5000\,\mathrm{t/d}$新型干法生产线窑尾配套的四电场电除尘器连续7天的运行数据，旨在通过分析入口条件、操作参数与出口粉尘浓度的复杂关系，设计典型工况划分方法，并在确保出口粉尘浓度满足国家超低排放标准（$\le 10\ \text{mg/Nm}^3$）\cite{mee2024opinion,mee2024qa}的前提下，为不同工况寻优出总电耗最低的操作参数组合（电压与振打周期），同时探究更严格排放标准下的能耗变化与应对策略，为水泥企业的绿色低碳生产提供技术支撑。

\section{问题分析}

\subsection{总体分析}

附件缺少二次电流、实际振打触发时刻、粉尘比电阻和粒径分布，又存在严重量程封顶。为避免“用复杂算法拟合无信息标签”，本文先回答\emph{数据能否识别}，再决定\emph{何处需要物理先验}。总体框架如图\ref{fig:framework}。

\begin{figure}[H]
\centering
\begin{tikzpicture}[
  node distance=8mm and 7mm,
  block/.style={rounded corners=1.4mm,draw=espblue,fill=esplight,minimum height=10mm,text width=3.05cm,align=center,font=\small},
  result/.style={rounded corners=1.4mm,draw=espteal,fill=black!2,minimum height=10mm,text width=3.05cm,align=center,font=\small},
  arr/.style={-{Stealth[length=2mm]},thick,draw=espgray}
]
\node[block] (raw) {分钟级原始数据\\10,080 行};
\node[block,right=of raw] (audit) {连续性、缺失、异常\\右删失与共线性审计};
\node[block,right=of audit] (evidence) {证据分层\\数据可识别/弱识别/先验主导};
\node[block,below=of audit] (regime) {15 min 稳健平滑\\GMM+Viterbi 工况识别};
\node[block,left=of regime] (power) {$U_i^2+1/T_i$\\能耗代理模型};
\node[block,right=of regime] (twin) {D--A 穿透骨架\\积灰+再飞扬数字孪生};
\node[result,below=of regime] (opt) {95\%机会约束\\差分进化+SLSQP};
\node[result,left=of opt] (q3) {两类工况策略\\边际收益与优先级};
\node[result,right=of opt] (q4) {$10\rightarrow5\mgNm$\\可行性与增量电耗};
\draw[arr] (raw)--(audit);
\draw[arr] (audit)--(evidence);
\draw[arr] (audit)--(regime);
\draw[arr] (audit)--(power);
\draw[arr] (evidence)--(twin);
\draw[arr] (regime)--(opt);
\draw[arr] (power)--(opt);
\draw[arr] (twin)--(opt);
\draw[arr] (opt)--(q3);
\draw[arr] (opt)--(q4);
\end{tikzpicture}
\caption{右删失信息下的协同优化建模框架}
\label{fig:framework}
\end{figure}

\subsection{各问题求解思路}

\begin{enumerate}
  \item \textbf{问题一：先判定数据能否支持关系识别。} 先检查时间连续性、缺失、边界堆积和闭环控制偏差；对疑似封顶读数采用删失似然，再用电除尘机理补充数据无法辨识的电压、温度、积灰和再飞扬方向。
  \item \textbf{问题二：把连续生产过程转化为少量可执行工况。} 以入口浓度和温度为主要工况轴，流量用于刻画类内负荷波动；通过时间正则化聚类得到稳定区间，再分别建立功率代理和排放风险模型，使用智能算法求解八个控制量的最低功率组合。
  \item \textbf{问题三：从最优点反推控制规律。} 选择入口质量负荷差异最大的两个稳定工况，通过历史均值对照、边际收益和活跃约束判断何时优先调电压、何时优先调振打周期。
  \item \textbf{问题四：先判可行，再谈增幅。} 在相同工况、风险水平和设备边界下比较 $10\mgNm$ 与 $5\mgNm$；若原边界不可行，则把“不可行”作为基础结论，只在明确的能力扩展情景中计算条件增幅。
\end{enumerate}

\section{模型假设与符号说明}

\subsection{基本假设}

\begin{enumerate}
  \item 入口浓度、流量及温度的分钟数据代表进入电除尘器的总体工况，时间戳无系统性错位；
  \item 根据连续变量在 50 处形成点质量且不存在更高读数的统计特征，本文把 $C_{\mathrm{out}}=50\mgNm$ 作为疑似上限删失事件，即在该工作假设下真实值不低于 50；另以先验压力测试评估假设风险，50 个出口缺失值不插补；
  \item 二次电流缺失时，总功率可由电压平方和振打频率的代理项描述；该关系仅用于历史安全范围附近的控制优化；
  \item 四场分场贡献、温度效应、积灰与再飞扬系数由公开机理研究提供先验，其不确定性通过蒙特卡洛传播，不解释为附件直接估计；
  \item 未知振打相位在工况内近似均匀，优化考虑相位积分后的分位风险；现场执行时四场振打必须错峰；
  \item 除问题 4 的条件情景外，电压与振打周期均不超出历史观测边界。
\end{enumerate}

\subsection{符号说明}

\begin{table}[H]
\centering
\caption{主要符号与单位}
\label{tab:symbols}
\begin{tabularx}{0.96\textwidth}{C{1.9cm}L C{2.2cm}L}
\toprule
符号 & 含义 & 符号 & 含义\\
\midrule
$C_{\mathrm{in}}$ & 入口粉尘浓度，$\gNm$ & $C_{\mathrm{out}}$ & 出口粉尘浓度，$\mgNm$\\
$Q$ & 烟气流量，$\NmH$ & $\Theta$ & 入口烟气温度，$^\circ\mathrm{C}$\\
$U_i$ & 第 $i$ 电场二次电压，$\kV$ & $T_i$ & 第 $i$ 电场振打周期，$\mathrm{s}$\\
$P$ & 总除尘功率，$\kW$ & $w_i$ & 第 $i$ 电场有效迁移贡献权重\\
$D$ & 等效去除指数 & $D_0$ & 基准工况等效去除指数\\
$L_t$ & 入口粉尘质量负荷相对值 & $M_{i,t}$ & 第 $i$ 电场极板积灰隐状态\\
$c$ & 仪表删失阈值，$50\mgNm$ & $\alpha$ & 机会约束置信水平，0.95\\
\bottomrule
\end{tabularx}
\end{table}
